% chapters/appendix-ch03-sql.tex
\chapter{SQL Appendix: Chapter 3 (Keys and Identity)}

\section{DDL}
\begin{verbatim}
CREATE TABLE student (
  student_id INT PRIMARY KEY AUTO_INCREMENT,
  full_name VARCHAR(120) NOT NULL,
  email VARCHAR(200) NULL UNIQUE
) ENGINE=InnoDB;

CREATE TABLE course (
  course_id INT PRIMARY KEY AUTO_INCREMENT,
  code VARCHAR(20) NOT NULL UNIQUE,
  title VARCHAR(200) NOT NULL
) ENGINE=InnoDB;

CREATE TABLE enrollment (
  student_id INT NOT NULL,
  course_id INT NOT NULL,
  term VARCHAR(10) NOT NULL,
  grade VARCHAR(2) NULL,
  PRIMARY KEY (student_id, course_id, term),
  CONSTRAINT fk_enroll_student FOREIGN KEY (student_id) REFERENCES student(student_id),
  CONSTRAINT fk_enroll_course FOREIGN KEY (course_id) REFERENCES course(course_id)
) ENGINE=InnoDB;
\end{verbatim}

\section{Sample data}
\begin{verbatim}
INSERT INTO student (full_name, email)
VALUES
  ('Marta Lima', 'marta@uni.edu'),
  ('Rui Nunes', NULL);

INSERT INTO course (code, title)
VALUES
  ('DB101', 'Database Design'),
  ('CS120', 'Programming Fundamentals');

INSERT INTO enrollment (student_id, course_id, term, grade)
VALUES
  (1, 1, '2025-FALL', 'A'),
  (1, 1, '2026-SPR', NULL),
  (2, 2, '2026-SPR', NULL);
\end{verbatim}

\section{Example queries}
\begin{verbatim}
-- Q1: Enrollments per student and term
SELECT student_id, course_id, term
FROM enrollment
WHERE student_id = 1;

-- Q2: Students with duplicate emails (should be none)
SELECT email, COUNT(*)
FROM student
WHERE email IS NOT NULL
GROUP BY email
HAVING COUNT(*) > 1;
\end{verbatim}
